\section{Complex heat time and real-time translations}
\label{sec:complex-heat}
The complex heat extension used in the boundary argument can be constructed
from the real Euclidean semigroup alone. This also fixes the sign of the
real-time translation group. Let $(E_t)_{t\geq0}$ be a strongly continuous
semigroup of self-adjoint contractions on a complex Hilbert space, with
$E_0=I$. Each $E_t$ is positive because $E_t=E_{t/2}^2$.

We use continuous functional calculus for a bounded self-adjoint operator:
continuous scalar functions on its spectrum act as a unital isometric
star homomorphism, with the usual composition rule. This is the
self-adjoint instance of the continuous calculus in \cite[Theorem 2.19]{Williams}.
All uses below are for a positive contraction and scalar functions continuous
at zero; an imaginary power at zero is never inserted into this calculus.

\begin{proposition}\label{prop:complex-heat}
There is a holomorphic operator-valued function $E(z)$ on
$\operatorname{Re}z>0$ such that
\[
 \|E(z)\|\leq1,\quad E(z+w)=E(z)E(w),\quad
 E(z)^*=E(\overline z),\quad E(t)=E_t\quad(t>0).
\]
Each $E_t$, $t>0$, is injective and has dense range. There is a strongly
continuous unitary group $V(t)$, $t\in\mathbb R$, characterized by
\begin{equation}\label{eq:complex-heat-edge}
 V(t)E(z)=E(z)V(t)=E(z-it)\qquad(\operatorname{Re}z>0).
\end{equation}
Every bounded operator commuting with all $E_t$ commutes with $E(z)$ and
$V(t)$. If $E_t\Omega=\Omega$ for all $t$, then
$E(z)\Omega=V(t)\Omega=\Omega$.
\end{proposition}
\begin{proof}
Put $A=E_1$. Its spectrum is contained in $[0,1]$: positivity excludes
negative spectrum and the contraction bound excludes values above one.
For $\operatorname{Re}z>0$, define
\[
 f_z(\lambda)=\begin{cases}\exp(z\log\lambda),&\lambda>0,\\0,&\lambda=0,\end{cases}
 \qquad E(z)=f_z(A).
\]
This scalar function is continuous at zero and has magnitude at most one.
The norm bound, product law, and adjoint identity follow from the calculus.
Holomorphy holds in the uniform norm on the scalar functions, hence in
operator norm. Indeed, on a compact neighborhood of any such $z$ choose
$\delta>0$ with $\operatorname{Re}z\geq2\delta$. The functions
$\lambda^\delta|\log\lambda|^j$ are bounded on $(0,1]$ and tend to zero at
zero for each fixed $j$. Taylor's formula for the scalar exponential then
gives, uniformly in $\lambda\in[0,1]$ and for sufficiently small $h$,
\[
 |f_{z+h}(\lambda)-f_z(\lambda)-h\lambda^z\log\lambda|
 \leq C_\delta |h|^2.
\]
At zero the derivative is defined to be zero. This proves complex
differentiability into the continuous-function space.

For a positive integer $n$, put $B=E_{1/n}$. The semigroup law gives
$B^n=A$. Functional calculus on the positive operator $B$ yields
$(B^n)^{1/n}=B$, because $(\lambda^n)^{1/n}=\lambda$ for
$\lambda\geq0$. The composition rule thus gives
$E(p/n)=A^{p/n}=B^p=E_{p/n}$ for positive integers $p,n$.
Approximate any positive real $t$ by positive rationals. Operator-norm
continuity of $E(z)$ and strong continuity of the original semigroup show,
after application to each vector, that $E(t)=E_t$.

If $E_tx=0$, then
$\|E_{t/2}x\|^2=\langle x,E_tx\rangle=0$.
Iteration gives $E_{t/2^n}x=0$ for every $n$, and strong continuity at zero
implies $x=0$. Self-adjointness then gives dense range, since the orthogonal
complement of the range is the kernel of the adjoint.
For real $t$ and $r>0$, the product and adjoint identities imply
\[
 E(r-it)^*E(r-it)=E(2r)=E_r^*E_r,
 \qquad \|E(r-it)x\|=\|E_rx\|.
\]
Define an isometry initially on the dense range of $A$ by
$V(t)Ax=E(1-it)x$. Injectivity of $A$ makes the definition unambiguous;
the displayed norm identity extends it uniquely to an isometry of the
whole Hilbert space.

For $\operatorname{Re}z>0$, both equalities in
\eqref{eq:complex-heat-edge} may be checked on $Ax$. For example,
\[
 V(t)E(z)Ax=V(t)AE(z)x=E(1-it)E(z)x
            =E(z-it)Ax.
\]
The other order is identical. Thus these are bounded-operator identities.
In particular, on $Ax$,
\[
 V(t)V(s)Ax=V(t)E(1-is)x=E(1-i(s+t))x=V(t+s)Ax.
\]
Density proves the group law. Since $V(0)=I$, each isometry has inverse
$V(-t)$ and is unitary. For $Ax$, the orbit is continuous by holomorphy
of $E(1-it)$; approximation by the dense range of $A$, with the unitary
norm bound on both sides, proves strong continuity for every vector.

Commutation with $A$ implies commutation with every continuous function of
$A$: approximate such a function uniformly on its real spectrum by
polynomials. Thus a bounded operator commuting with all $E_t$ commutes
with $E(z)$. Checking on $Ax$ then gives commutation with $V(t)$.
Finally $A\Omega=\Omega$ implies $f(A)\Omega=f(1)\Omega$, first for
polynomials and then for continuous $f$. Apply this to $f_z$ and then the
definition of $V(t)$ to obtain the vacuum assertions.
\end{proof}

Suppose now that spatial translations $U(x)$ form a strongly continuous
unitary representation commuting with the heat semigroup and fixing
$\Omega_E$. The proposition
defines a strongly continuous unitary representation of all real translations,
\[
 T_E(a^0,\mathbf a)=V(a^0)U(\mathbf a).
\]
The group law follows from commutation. Joint strong continuity follows by
adding and subtracting $V(a^0)U(\mathbf b)\xi$ and using the unitary norms.
Moreover $E(\varepsilon-it)\xi=V(t)E_\varepsilon\xi\to V(t)\xi$ as
$\varepsilon\downarrow0$, which justifies the boundary notation $E(-it)$
used in the heat-vector chapter. No operator-norm continuity at the boundary
is required.

\begin{lemma}\label{lem:time-fixed-heat}
The vectors fixed by every $V(t)$ are exactly the vectors fixed by every
$E_s$. Consequently, if the selected heat ground space is
$\mathbb C\Omega_E$, the same is true of the subspace fixed by all real
translations.
\end{lemma}
\begin{proof}
A heat-fixed vector satisfies $A\xi=\xi$, so functional calculus and the
definition of $V$ give $V(t)\xi=\xi$. Conversely suppose $V(t)\xi=\xi$.
The holomorphic vector function $g(z)=E(z)\xi$ satisfies
$g(z-it)=E(z)V(t)\xi=g(z)$. Its derivative in the imaginary direction is
zero; complex differentiability therefore makes its full derivative zero.
It is constant on the connected right half-plane. Since
$E_s\xi\to\xi$ as $s\downarrow0$, this constant is $\xi$.
A vector invariant under all translations is, in particular, invariant under
$V(t)$, and hence belongs to the heat ground space. The vacuum itself is
translation-invariant, which proves the last assertion.
\end{proof}

Combining this lemma with the isometric translation intertwiner in
Proposition~\ref{prop:word-motion} proves uniqueness of the invariant vacuum
in the reconstructed field Hilbert space: an invariant $\xi$ maps to
$c\Omega_E=\mathcal V(c\Omega)$, and injectivity of the intertwiner
$\mathcal V$ gives $\xi=c\Omega$.
